\documentclass[11pt,letterpaper,oneside]{article}

\usepackage[utf8]{inputenc}
\usepackage[T1]{fontenc}
\usepackage[provide=*,spanish]{babel}
\usepackage{lmodern}
\usepackage{geometry}
\usepackage{xcolor}
\usepackage{array}
\usepackage{booktabs}
\usepackage{tabularx}
\usepackage{longtable}
\usepackage{verbatim}
\usepackage{url}
\usepackage{hyperref}

\geometry{top=2.1cm,bottom=2.2cm,left=2.3cm,right=2.3cm}
\setlength{\parindent}{0pt}
\setlength{\parskip}{0.55em}
\setlength{\emergencystretch}{3em}
\renewcommand{\arraystretch}{1.18}
\raggedbottom

\definecolor{uvgblue}{HTML}{123B65}
\definecolor{purple}{HTML}{6842D8}
\definecolor{cyan}{HTML}{087F8C}
\definecolor{softblue}{HTML}{EAF3FA}
\definecolor{softgray}{HTML}{F3F5F7}

\hypersetup{
  colorlinks=true,
  linkcolor=uvgblue,
  urlcolor=cyan,
  pdftitle={Guion de exposición -- Laboratorio 01},
  pdfauthor={DragonsSlayers 2.0}
}

\newcommand{\codigo}[1]{\texttt{#1}}
\newcommand{\checkbox}{\fbox{\rule{0pt}{1.1ex}\hspace{1.1ex}}}
\newcommand{\accion}[1]{%
  \par\noindent
  \colorbox{softblue}{\parbox{0.94\linewidth}{\textbf{En pantalla:} #1}}
  \par
}
\newcommand{\nota}[1]{%
  \par\noindent
  \colorbox{softgray}{\parbox{0.94\linewidth}{\textbf{Nota de grabación:} #1}}
  \par
}
\newcommand{\habla}[2]{%
  \par\noindent
  {\color{purple}\textbf{#1:}} #2
  \par
}
\newcommand{\escena}[3]{%
  \subsection*{#1 \hfill \normalfont\color{cyan}#2}
  \phantomsection
  \addcontentsline{toc}{subsection}{#1 -- #2}
  \textbf{Responsable:} #3
}

\begin{document}

\hypersetup{pageanchor=false}

\begin{titlepage}
  \centering
  \vspace*{1cm}
  {\Large\bfseries Universidad del Valle de Guatemala\par}
  \vspace{0.25cm}
  {\large Construcción de Compiladores -- CC-3032, sección 10\par}
  \vfill
  \color{uvgblue}\rule{0.88\textwidth}{1.2pt}\color{black}
  \vspace{0.7cm}
  {\Huge\bfseries Guion de exposición en video\par}
  \vspace{0.35cm}
  {\LARGE Laboratorio 01\par}
  \vspace{0.25cm}
  {\Large Uso de generadores de analizadores léxicos y sintácticos\par}
  \vspace{0.65cm}
  {\Huge\color{purple}\bfseries Analizador de Compiscript\par}
  \vspace{0.7cm}
  \color{uvgblue}\rule{0.88\textwidth}{1.2pt}\color{black}
  \vfill
  \begin{tabular}{>{\bfseries}rl}
    Grupo: & DragonsSlayers 2.0 \\
    Integrante 1: & Hugo Daniel Barillas Ajín -- 23556 \\
    Integrante 2: & Ernesto Ascencio -- 23009 \\
    Duración objetivo: & 9 minutos 40 segundos \\
    Restricción: & máximo 10 minutos \\
  \end{tabular}
  \vfill
  {\large Guatemala, 2026\par}
\end{titlepage}

\hypersetup{pageanchor=true}
\pagenumbering{roman}

\section*{Cómo utilizar este guion}
\phantomsection
\addcontentsline{toc}{section}{Cómo utilizar este guion}

Este guion corresponde a una exposición continua de 9 minutos 40 segundos,
distribuida únicamente entre Hugo y Ernesto. Cada escena indica el tiempo objetivo,
la acción que debe mostrarse y un texto que puede leerse casi literalmente. Las cajas
azules no se leen; describen la evidencia visual.

Los tiempos dejan 20 segundos de margen respecto del máximo. Los ocho archivos deben
estar abiertos o accesibles antes de grabar. No se debe afirmar que Compiscript se
ejecuta o se compila a código objeto: el alcance es exclusivamente léxico y
sintáctico.

\tableofcontents
\clearpage
\pagenumbering{arabic}

\section{Preparación antes de grabar}

\subsection{Ventanas y archivos}

\begin{enumerate}
  \item Aplicación web abierta en \url{http://localhost:3000}.
  \item Editor con \codigo{Compiscript.g4}, \codigo{analyze.ts} y
        \codigo{antlrErrors.ts}.
  \item Terminal ubicada en la raíz del repositorio.
  \item Explorador en \path{examples/rubric/low} y
        \path{examples/rubric/medium}.
  \item Opcional: aplicación portable de Electron ya construida.
\end{enumerate}

Ejecutar antes de grabar:

\begin{verbatim}
npm install
npm run generate
npm run check
npm test
npm run dev
\end{verbatim}

La suite debe mostrar \textbf{3 archivos aprobados y 24 pruebas aprobadas}. Aumentar
el tamaño de letra del navegador y de la terminal para que tokens, líneas, columnas y
mensajes sean legibles.

\subsection{Material obligatorio}

\begin{itemize}
  \item Los ocho casos de \path{examples/rubric}.
  \item Los resúmenes \codigo{0/0}, \codigo{3/0}, \codigo{0/3} y
        \codigo{2/2} para baja y media.
  \item Al menos un diagnóstico con fase, línea, columna, símbolo y mensaje.
  \item Tabla de tokens y árbol de parseo.
  \item Carga de archivos \codigo{.cps}, exportaciones, CLI y pruebas.
  \item Mención breve de ANTLR, Vite, React y Electron.
\end{itemize}

\section{Distribución definitiva del tiempo}

\begin{longtable}{c p{0.48\textwidth} p{0.22\textwidth}}
\toprule
Tiempo & Tema & Responsable \\
\midrule
\endfirsthead
\toprule
Tiempo & Tema & Responsable \\
\midrule
\endhead
00:00--00:35 & Presentación y objetivo & Hugo \\
00:35--01:15 & Teoría léxica y sintáctica & Ernesto \\
01:15--01:55 & ANTLR, arquitectura y alcance & Hugo \\
01:55--02:30 & Interfaz y definición de complejidades & Hugo \\
02:30--03:00 & Caso válido de complejidad baja & Hugo \\
03:00--04:00 & Tres casos con errores de complejidad baja & Hugo \\
04:00--04:40 & Caso válido de complejidad media & Ernesto \\
04:40--05:50 & Tres casos con errores de complejidad media & Ernesto \\
05:50--06:40 & Recuperación, ubicación y mensajes & Ernesto \\
06:40--07:20 & Tokens, árbol y exportaciones & Hugo \\
07:20--08:00 & CLI y 24 pruebas automatizadas & Ernesto \\
08:00--08:40 & Ejecución web y aplicación de escritorio & Hugo \\
08:40--09:15 & Comprobación de la rúbrica & Ernesto y Hugo \\
09:15--09:40 & Conclusiones y cierre & Ernesto y Hugo \\
\bottomrule
\end{longtable}

\nota{No agregar improvisaciones sin recortar otra escena. Mostrar cada resumen antes
de cambiar de archivo, pero no recorrer completa la tabla de tokens ocho veces.}

\section{Guion por escenas}

\escena{Escena 1. Presentación}{00:00--00:35}{Hugo}

\accion{Mostrar portada con nombre del laboratorio, ANTLR 4, TypeScript, Hugo y
Ernesto.}

\habla{Hugo}{Buenos días. Somos Hugo Daniel Barillas Ajín y Ernesto Ascencio, del
grupo DragonsSlayers 2.0. Presentamos un analizador léxico y sintáctico para archivos
Compiscript con extensión punto cps. El objetivo es reconocer tokens, validar la
estructura del programa y reportar múltiples errores comprensibles sin detenerse en
el primero.}

\escena{Escena 2. Fundamento teórico}{00:35--01:15}{Ernesto}

\accion{Mostrar un diagrama breve: caracteres, lexer, tokens, parser y árbol.}

\habla{Ernesto}{El lexer transforma caracteres en tokens como palabras reservadas,
identificadores, literales y operadores. Un símbolo que no pertenece al vocabulario
produce un error léxico. El parser consume esos tokens y verifica las producciones de
la gramática; por ejemplo, una declaración debe terminar con punto y coma. Una
omisión estructural es un error sintáctico. El resultado correcto del parser es un
árbol de parseo que inicia en program y alcanza EOF.}

\escena{Escena 3. ANTLR, arquitectura y alcance}{01:15--01:55}{Hugo}

\accion{Mostrar \codigo{Compiscript.g4}, \codigo{src/generated} y un diagrama con
\codigo{analyzeInput} conectado a React, CLI y Electron.}

\habla{Hugo}{ANTLR 4 genera el lexer, parser, listener y visitor TypeScript desde
Compiscript punto g cuatro. La gramática cubre tipos, arreglos, funciones, clases,
expresiones, condicionales y ciclos. El núcleo analyzeInput es compartido por la web,
la consola y Electron, por lo que todas las vistas usan el mismo resultado. El
alcance termina en análisis léxico y sintáctico; no se ejecuta Compiscript ni se
realiza comprobación semántica.}

\escena{Escena 4. Interfaz y criterios de complejidad}{01:55--02:30}{Hugo}

\accion{Mostrar la interfaz y las carpetas de complejidad baja y media.}

\habla{Hugo}{La interfaz permite análisis completo, solo lexer, edición y carga de
archivos punto cps. La rúbrica pide cuatro casos bajos y cuatro medios. Todo caso bajo
conserva tres tipos de variables, una constante, dos operadores aritméticos, un if,
un while y un foreach. Los casos medios agregan un arreglo usado, dos clases, dos
objetos, dos funciones y dos llamadas.}

\escena{Escena 5. Complejidad baja válida}{02:30--03:00}{Hugo}

\accion{Cargar \path{examples/rubric/low/valid.cps}, señalar brevemente sus
construcciones y analizarlo.}

\habla{Hugo}{El primer archivo cumple toda la complejidad baja. El analizador recorre
112 tokens y lo acepta con cero errores léxicos y cero sintácticos. Esto demuestra el
primer criterio y permite observar el árbol completo.}

\escena{Escena 6. Errores en complejidad baja}{03:00--04:00}{Hugo}

\accion{Cargar los otros tres archivos de \codigo{low}. Dejar visible cada resumen:
\codigo{3/0}, \codigo{0/3} y \codigo{2/2}.}

\habla{Hugo}{El segundo archivo conserva la misma complejidad e introduce arroba,
numeral y dólar: reporta tres errores léxicos y ningún error sintáctico. El tercero
omite tres puntos y coma: reporta cero errores léxicos y tres sintácticos. El cuarto
combina dos símbolos inválidos con dos omisiones y produce exactamente dos errores
de cada fase. En todos los casos el procesamiento continúa hasta las sentencias
finales.}

\escena{Escena 7. Complejidad media válida}{04:00--04:40}{Ernesto}

\accion{Cargar \path{examples/rubric/medium/valid.cps}; señalar el arreglo, las
clases Estudiante y Profesor, las dos instancias, las funciones y sus llamadas.}

\habla{Ernesto}{El caso medio mantiene todos los requisitos bajos y agrega un arreglo
de notas utilizado por for y foreach, dos clases, dos objetos creados con new, las
funciones sumar y mostrarCurso, y sus llamadas. Se reconocen 263 tokens y el archivo
es aceptado con cero errores en ambas fases.}

\escena{Escena 8. Errores en complejidad media}{04:40--05:50}{Ernesto}

\accion{Cargar los otros tres archivos de \codigo{medium}. Mostrar en cada resultado
los conteos \codigo{3/0}, \codigo{0/3} y \codigo{2/2}.}

\habla{Ernesto}{El caso léxico medio contiene tres símbolos no reconocidos y conserva
cero errores del parser. El caso sintáctico omite tres puntos y coma y conserva cero
errores del lexer. El caso mixto reporta dos errores léxicos y dos sintácticos. Las
construcciones de complejidad media permanecen en los tres archivos, de modo que no
son ejemplos mínimos desconectados de la rúbrica.}

\escena{Escena 9. Recuperación y diagnósticos}{05:50--06:40}{Ernesto}

\accion{Ampliar dos errores del caso mixto y señalar fase, línea, columna, símbolo y
mensaje. Mostrar \codigo{tokenStream.fill()} y \codigo{DefaultErrorStrategy}.}

\habla{Ernesto}{Fill obliga al lexer a recorrer la entrada completa. Después,
DefaultErrorStrategy permite al parser insertar o descartar tokens conceptualmente y
sincronizarse para continuar. Cada diagnóstico indica si proviene del lexer o parser,
la línea y columna base uno, el símbolo relacionado y una explicación en español.
Además se agrupan caracteres inválidos contiguos, se eliminan duplicados y se filtran
consecuencias inmediatas para evitar cascadas repetidas.}

\escena{Escena 10. Tokens, árbol y exportaciones}{06:40--07:20}{Hugo}

\accion{Volver al caso medio válido. Mostrar tabla de tokens, árbol visual y textual,
y los botones de descarga.}

\habla{Hugo}{La tabla conserva índice, tipo, lexema, línea, columna y canal de cada
token. El árbol puede explorarse visualmente o descargarse como texto. También se
exportan la gramática, la entrada, tokens CSV o JSON y el resultado completo. Todo se
genera localmente en el navegador, sin enviar el código a un servidor.}

\escena{Escena 11. CLI y pruebas}{07:20--08:00}{Ernesto}

\accion{Ejecutar \codigo{npm test} y mostrar 3 archivos y 24 pruebas aprobadas. Luego
mostrar un comando CLI de la carpeta \codigo{rubric}.}

\habla{Ernesto}{Vitest ejecuta 24 pruebas. Ocho leen directamente la matriz de la
rúbrica y verifican aceptación, conteos y construcciones de complejidad. La CLI
reutiliza el mismo motor: un archivo válido termina con código cero y uno inválido con
código uno, que representa rechazo esperado y no una excepción.}

\escena{Escena 12. Ejecución y escritorio}{08:00--08:40}{Hugo}

\accion{Mostrar \codigo{npm run dev}, la aplicación web y, sin compilar durante el
video, el ejecutable o la configuración de Electron.}

\habla{Hugo}{Durante desarrollo se instala con npm install y se inicia con npm run
dev. La versión de escritorio usa Electron y puede construirse como portable o como
instalador NSIS para Windows x64. El ejecutable final incluye la aplicación, por lo
que el usuario no necesita instalar Node, npm ni Java.}

\escena{Escena 13. Comprobación de la rúbrica}{08:40--09:15}{Ernesto y Hugo}

\accion{Mostrar una tabla de ocho filas con los resultados verificados y una captura
de un diagnóstico con ubicación.}

\habla{Ernesto}{Quedan demostrados los cuatro criterios de complejidad baja y los
cuatro de complejidad media: válido, léxico, sintáctico y mixto.}

\habla{Hugo}{También se cumple la ubicación correcta por línea y columna, y el reporte
inteligible en español con fase y símbolo o token relacionado. Así se cubren los diez
criterios de evaluación con archivos reproducibles.}

\escena{Escena 14. Conclusiones}{09:15--09:40}{Ernesto y Hugo}

\accion{Volver a la portada final con las palabras gramática, generación,
recuperación y validación.}

\habla{Ernesto}{ANTLR permitió definir el lenguaje declarativamente y recuperar
múltiples errores en una sola ejecución.}

\habla{Hugo}{La matriz y sus pruebas conectan cada requisito con evidencia concreta.
El resultado es un analizador reutilizable, documentado y verificable. Muchas
gracias.}

\section{Referencia para la demostración}

\subsection{Matriz oficial de la rúbrica}

Hugo presenta los cuatro casos de \codigo{low}; Ernesto presenta los cuatro de
\codigo{medium}.

\begin{center}
\begin{tabular}{l r r r}
\toprule
Caso & Tokens & LEX & SYN \\
\midrule
\path{low/valid.cps} & 112 & 0 & 0 \\
\path{low/lexer_errors.cps} & 112 & 3 & 0 \\
\path{low/parser_errors.cps} & 109 & 0 & 3 \\
\path{low/lexer_parser_errors.cps} & 110 & 2 & 2 \\
\path{medium/valid.cps} & 263 & 0 & 0 \\
\path{medium/lexer_errors.cps} & 263 & 3 & 0 \\
\path{medium/parser_errors.cps} & 260 & 0 & 3 \\
\path{medium/lexer_parser_errors.cps} & 261 & 2 & 2 \\
\bottomrule
\end{tabular}
\end{center}

\subsection{Correspondencia de complejidad}

\begin{tabularx}{\textwidth}{>{\bfseries}p{0.24\textwidth} X}
\toprule
Nivel & Construcciones conservadas en sus cuatro archivos \\
\midrule
Bajo & Variables \codigo{integer}, \codigo{string} y \codigo{boolean}; constante;
suma y multiplicación; \codigo{if}; \codigo{while}; \codigo{foreach}. \\
Medio & Todo lo anterior; arreglo declarado y usado; clases Estudiante y Profesor;
dos expresiones \codigo{new}; funciones sumar y mostrarCurso; dos llamadas. \\
\bottomrule
\end{tabularx}

\subsection{Comandos rápidos}

\begin{verbatim}
npm run cli -- examples/rubric/low/valid.cps
npm run cli -- examples/rubric/low/lexer_errors.cps
npm run cli -- examples/rubric/low/parser_errors.cps
npm run cli -- examples/rubric/low/lexer_parser_errors.cps
npm run cli -- examples/rubric/medium/valid.cps
npm run cli -- examples/rubric/medium/lexer_errors.cps
npm run cli -- examples/rubric/medium/parser_errors.cps
npm run cli -- examples/rubric/medium/lexer_parser_errors.cps
\end{verbatim}

Los dos casos válidos terminan con código 0. Los seis casos incorrectos terminan con
código 1 porque el rechazo es intencional.

\subsection{Carpeta complementaria de recuperación}

\path{examples/error_distribution} conserva tres archivos adicionales para estudiar
agrupación, recuperación y alternancia prolongada:

\begin{center}
\begin{tabular}{l r r r}
\toprule
Archivo & Tokens & LEX & SYN \\
\midrule
\path{lexer_errors.cps} & 39 & 3 & 0 \\
\path{parser_errors.cps} & 64 & 0 & 7 \\
\path{lexer_parser_alternating_errors.cps} & 72 & 4 & 4 \\
\bottomrule
\end{tabular}
\end{center}

Estos archivos son evidencia suplementaria; durante el video tienen prioridad los
ocho casos que corresponden directamente a la rúbrica.

\section{Lista final de verificación}

\begin{tabularx}{\textwidth}{p{0.05\textwidth} X}
\checkbox & Duración ensayada menor o igual a 9:40. \\
\checkbox & Hablan únicamente Hugo y Ernesto. \\
\checkbox & Se muestran los ocho archivos de la matriz. \\
\checkbox & Se ven los conteos 0/0, 3/0, 0/3 y 2/2 en ambos niveles. \\
\checkbox & Se muestra línea, columna, fase, símbolo y mensaje en español. \\
\checkbox & Se explica la diferencia entre lexer y parser. \\
\checkbox & Se muestran tokens y árbol de parseo. \\
\checkbox & Se ejecuta \codigo{npm test}: 3 archivos y 24 pruebas aprobadas. \\
\checkbox & Se menciona el alcance sin semántica ni ejecución de Compiscript. \\
\checkbox & El enlace final de YouTube se coloca en el README. \\
\end{tabularx}

\vfill
\begin{center}
\color{uvgblue}\rule{0.75\textwidth}{0.8pt}\color{black}

\textbf{DragonsSlayers 2.0}\par
Laboratorio 01 -- Construcción de Compiladores CC-3032\par
Guion definitivo: 9 minutos 40 segundos.
\end{center}

\end{document}
